\section{Compressed connectivity of cut-open normal surfaces}
\label{sec:cut-complement}

The hierarchy needs more than a component census of its cutting surface.
It must also locate the connected pieces of the ambient manifold after the
cut. Expanding a binary normal vector into individual discs loses the desired
complexity bound. This continuation constructs a linear-size interval system
for the connectivity of the cut-open manifold and authenticates its count
with the existing independent orbit-trace checker. It is a native geometric
operation toward compressed cutting, not a full implementation of that
hierarchy obligation.

The current source contract remains the existing validated finite connected
compact orientable triangulation with one torus boundary. The surface is
supplied by admissible nonnegative integral normal coordinates. It need not
be connected, primitive, two-sided or incompressible. No correspondence with
an input knot diagram is inferred. The result counts the components of
$M\setminus\operatorname{int}\nu(F)$, where $\nu(F)$ is a normal regular
neighborhood of the surface, including the one-sided case.

\subsection{Local chambers and affine face regions}

In a tetrahedron let $t_v$ be its four triangle counts and let $q$ be its
single occupied quadrilateral count, zero when no quadrilateral is present.
Cutting the tetrahedron along these normal discs gives
\[
                 1+q+\sum_{v=0}^{3}t_v
\]
connected chambers. Index the $q+1$ successive chambers of the quadrilateral
stack first, followed by four arms of $t_v$ chambers each, numbered outward
from vertex $v$. The triangle arms meet the appropriate endpoint chamber
of the quadrilateral stack. With $q=0$, all four arms meet the central chamber.
This is a description of the chambers; adjacent chambers separated by the
surface are never identified by a local tree edge.

Quadrilateral type $j$ separates the vertex pairs
$A=\{0,j+1\}$ and $B=\{0,1,2,3\}\setminus A$.
Take the chain index to increase from the $A$ side to the $B$ side. On a
face, each normal arc cuts off one of its three vertices. The face therefore
has a central region and three ordered arms of arc regions. At its corner
$v$, the arm has length $t_v+q_{f,v}$, where $q_{f,v}$ is the count of
quadrilaterals whose face arc cuts off $v$.

The first $t_v$ face-arm regions map in order to the triangle chambers.
Only the face's singleton vertex on one side of the quadrilateral partition
has further regions. Its $q$ remaining regions map to chain indices
$0,\ldots,q-1$ on the $A$ side, and $q,\ldots,1$ on the $B$ side.
The central face region maps to the opposite chain endpoint. These are
at most two affine pieces per face arm, with slopes $+1$ or $-1$.

For a glued face, normal matching equates the two arc-arm lengths.
Its vertex permutation matches regions of the same rank from the corresponding
vertices, including its central region. Intersect the two arm partitions
at their triangle/quad breakpoints. There are at most three nonempty
rank intervals per arm. Composing the two affine chamber maps produces
an interval isometry between chamber indices. Add the single central-region
pair. Self-face gluings and order reversals are included; boundary faces
produce no identification.

\begin{proposition}[Binary cut-component reduction]
For $T$ source tetrahedra and maximum coordinate bit length $B$, the described
system has
\[
 N=T+D,\qquad K\le20T,\qquad
 \log(N+1)=O(B+\log(T+1)),
\]
where $D$ is the number of represented normal discs. Its interval orbits
are exactly the connected components of the cut-open manifold. Construction,
orbit counting and independent trace replay have polynomial bit complexity
in the source size and coordinate bit length; replay also charges its supplied
trace length. No represented normal disc or chamber family is expanded.
\end{proposition}
\begin{proof}
The normal discs form a quadrilateral stack with four compatible nested
triangle stacks near the vertices. Every disc separates one local chamber,
so the stated chain and arms account for all $1+q+\sum t_v$ local pieces.
Each chamber is connected. Cutting a face along its arcs gives precisely the
central region and its three corner arms. The rank descriptions above record
which chamber meets every such region. Matching and the face permutation
give every actual identification across a glued face, without adding an
identification across the cutting surface.

A path in the cut-open manifold can be perturbed to cross source faces
through their region interiors, avoiding their edges and vertices. Its
sequence of local chambers therefore gives an interval-orbit path.
Conversely every interval arrow is a genuine shared face region and joins
two connected local chambers. Chaining these paths proves the equivalence.
A one-sided surface merely causes its two local sides to be identified by
source face gluings; no global coorientation is assumed.

Summing the local counts gives $N=T+D$. Each glued face contributes at most
ten interval pairings, and there are at most $2T$ glued face pairs. Also
$D\le7T\,2^B$, proving the digit bound. Only a constant number of breakpoints
per face are sorted and composed, so interval construction has linear scalar
work in $T$, charged at its endpoint bit sizes. The established
Agol--Hass--Thurston interval algorithm gives polynomial orbit work in
$K$ and $\log(N+1)$\cite{aht-orbits}; the retained independent local trace
checker charges all supplied events. Finite manifold validation and exact
matching validation retain their existing polynomial costs.
\end{proof}

This is an elementary chamber-connectivity reduction using the existing
orbit machinery; no new literature-priority claim is made. Regina's
\code{cutAlong} supplies an independent small-input geometric oracle; its
API explicitly warns that the expanded triangulation can be very large
\cite{regina-cut-along}. The native API supplies the component count and
its source-bound proof, not that expanded triangulation.

\subsection{Independent geometry and binary multiplicity}

The producer numbers chain and triangle-arm intervals and composes face
maps by evaluating their endpoints. The independent checker derives each
face's singleton vertex from its quad partition, constructs its affine
segments and intersects their domains. It reconstructs the canonical chamber
pairings from the actual validated source and coordinates. It imports neither
the producer's chamber constructor nor the orbit planner.

The schema \code{normal-cut-components-v1} binds the normalized source
and coordinates, the cut-component count and the complete ordinary orbit
certificate. All pairings, the universe size and the final count must match
the independent construction and replay. Signed binary integer transport
is retained. Invalid schemas, changed sources or coordinates, changed counts
and altered traces do not authenticate a result. Cooperative cancellation
propagates. A capped orbit query contains neither a component count nor
a certificate, and a replay operation limit means unverified rather than
any topological conclusion.

Multiplicity must not be divided out as it is for some surface-census queries.
For $m\ge1$ parallel meridian discs in a solid torus the cut has $m$ pieces;
$m$ parallel vertex-link discs give $m+1$. For the retained one-tetrahedron
M\"obius-band vector, its $m$-fold normal vector gives
$\lfloor m/2\rfloor+1$ pieces. Odd multiplicity includes a one-sided sheet,
whereas even multiplicity consists of annular doubles. The chamber model
handles these cases directly on binary counts without producing their copies.

\subsection{Scope still missing for the hierarchy}

The new API is \code{normal\_complement\_components}.

Independent replay uses \code{verify\_normal\_complement\_certificate}.

It is not called by
the complete recognizer and supplies no knot verdict or new recognition
asymptotic bound. A cut can have exponentially many components; a count does
not list them or extract a usable compressed manifold representation.

Post-cut cells with full boundary-pattern transport, component descriptions,
regluing maps and provenance for lifting later discs are still absent.
The existing one-torus validator also does not validate arbitrary later
hierarchy pieces with multiple boundary components. Pattern-sensitive
surfaces, bounded rebuilds and logarithmic hierarchy depth remain unproved.
Thus this implements a concrete polynomial connectivity operation inside the
missing cutting requirement, without claiming that requirement complete.

\subsection{Geometric cross-checks and component-count measurements}
All 1,548 maintained tests pass in 288.894 seconds; the eight focused tests
pass in 0.284 seconds. They check actual-source chamber pairings, empty cuts,
parallel meridians and vertex links, one-sided sheets and their doubles,
mixed link/quad surfaces, source and trace mutations, producer-disabled
replay, operation/cycle caps and cancellation. A 16,385-bit coordinate
certificate passes ordinary binary JSON transport and independent replay.

The initial geometric pilot matches 1,635 standard-vertex or empty surfaces
on 48 source triangulations against Regina's cut-open component counts.
The integrated 10.894-second audit repeats those actual cuts with the native
source-bound proof and independent chamber reconstruction, and adds 384
compatible mixed-coordinate surfaces. Every one of the 2,019 expanded-cut
checks agrees. These are surface instances, not 2,019 distinct triangulations.
All native proofs replay with chamber and orbit producers disabled.

Thirteen further binary cases include parallel meridians on one-, four- and
sixteen-tetrahedron layered solid tori, with multiplicity one, $2^{128}$ or
$2^{500}$; a one-tetrahedron $2^{16384}$ meridian case; and large vertex-link
and odd/even M\"obius-sheet families. Their counts agree with the geometric
family formulas above, and their exact certificates independently replay.
No expanded Regina cut is attempted for the enormous multiplicities.
The sixteen-tetrahedron meridian family takes 68 orbit cycles for all three
scales; the one-tetrahedron meridian family takes four. These observations
support the implementation checks, while the general polynomial guarantee
comes from the proved chamber reduction and interval algorithm.

The audit retains 157 representative original-source proofs, including all
thirteen large cases. The publication review independently replays these,
checks the complete oracle records and authenticates 658 source/evidence
hashes together with the exact pilot and source-bank hashes.

Five paired timing rounds measure 80 completed component-count calls and
sixteen warmups, with identical-baseline A/A arms. Each input is a binary
multiple of the primitive meridian of the same one-tetrahedron source.
The baseline constructs a Regina cut-open triangulation and counts its
components, including its source/surface construction and result serialization.
The native arm constructs chamber intervals, counts their orbits, independently
replays the source-bound certificate and serializes its result. Thus the
comparison has the same requested count but different intermediate artifacts:
only the baseline constructs the expanded cut triangulation. It is not a
comparison of two complete cut-manifold constructors or two knot recognizers.
\begin{center}
\begin{tabular}{rrrrrrr}
\toprule
copies & cut $T$ & explicit ms & native ms & old/new & old A/A & new A/A\\
\midrule
1 & 31 & 0.151 & 0.414 & 0.353 & 0.914 & 1.023\\
16 & 226 & 0.404 & 0.456 & 0.935 & 1.020 & 1.056\\
128 & 1682 & 2.422 & 0.499 & 4.826 & 0.999 & 1.119\\
1024 & 13330 & 34.197 & 0.468 & 74.860 & 0.997 & 1.088\\
\bottomrule
\end{tabular}
\end{center}


At 1,024 copies the paired explicit/native ratio is 74.860, with median
34.197 versus 0.468 milliseconds. The explicit cut has 13,330 tetrahedra;
the compressed model has 3,073 chamber points, three interval pairings and
four orbit cycles. At 128 copies the ratio is 4.826. At one and sixteen
copies the ratios are 0.353 and 0.935: native Python validation and independent
proof replay have a larger small-input constant than the C++ expanded baseline.
These regressions are retained rather than presented as improvements.

The new A/A ratios range from 1.023 to 1.119 and the old ones from 0.914
to 1.020. Submillisecond observations fluctuate; the large ratio is not a
precise universal speedup. The important asymptotic distinction is the
binary count operation versus constructing every represented chamber.
No general recognition speedup or changed recognizer default follows.

Runtime release \code{cbbd59371}, the passing tests, both source audits and
the complete count-task measurements are retained.

The driver is \code{normal\_orbit\_research.cut\_complement}, with audit
and benchmark modes. Records are under \code{data/cut-complement-*} and
\code{data/complement-chamber-pilot.*}. The operator completes the local
cut-connectivity calculation under its existing source contract. It does
not extract or validate later hierarchy pieces, transport boundary patterns
or prove logarithmic depth. The full compressed-cutting obligation and
general quasi-polynomial recognition bound remain open.

